\section*{Chapitre \ref{ch:g3:sorting-living-things} --- Classer les êtres vivants}

\begin{solution}{exo:g3:sorting-living-things:1}
Un \omterm{def:g3:sorting-living-things:criterion}{critère} est la règle qui sert au tri. Bon~: un caractère du corps
que l'on peut observer --- «~a des \omterm{def:g1:animals-around-us:covering}{plumes}~», «~a six pattes~».
Évité~: le lieu où il vit ou ce qu'il est en train de faire --- «~vit
près de l'eau~», «~vole~».
\end{solution}

\begin{solution}{exo:g3:sorting-living-things:2}
\omterm{prop:g3:sorting-living-things:groups}{Insectes}~: six pattes et un corps en trois parties. \omterm{prop:g3:sorting-living-things:groups}{Poissons}~: des
nageoires et des \omterm{def:g1:animals-around-us:covering}{écailles}, une vie dans l'eau. \omterm{prop:g3:sorting-living-things:groups}{Oiseaux}~: des \omterm{def:g1:animals-around-us:covering}{plumes},
des ailes et un bec. \omterm{prop:g3:sorting-living-things:groups}{Mammifères}~: des poils, et des mères qui
nourrissent leurs petits de \omterm{def:g2:animals-grow-up:born}{lait}.
\end{solution}

\begin{solution}{exo:g3:sorting-living-things:3}
Abeille~: insecte. Poisson rouge~: poisson. Autruche~: oiseau.
Vache~: mammifère. Grenouille~: amphibien.
\end{solution}

\begin{solution}{exo:g3:sorting-living-things:4}
Elle a des poils et nourrit ses petits de \omterm{def:g2:animals-grow-up:born}{lait}, et elle n'a ni
\omterm{def:g1:animals-around-us:covering}{plumes} ni bec. Voler est quelque chose qu'elle fait, pas ce qui
définit un oiseau --- les caractères disent mammifère.
\end{solution}

\begin{solution}{exo:g3:sorting-living-things:5}
Elle n'a pas d'\omterm{def:g1:animals-around-us:covering}{écailles} ni le plan du corps à branchies et
nageoires d'un poisson --- et elle nourrit son baleineau de \omterm{def:g2:animals-grow-up:born}{lait}.
Nager est son habitude~; ses caractères disent mammifère.
\end{solution}

\begin{solution}{exo:g3:sorting-living-things:6}
Non. Le \omterm{def:g3:sorting-living-things:criterion}{critère} des \omterm{prop:g3:sorting-living-things:groups}{insectes} est six pattes (et un corps en trois
parties)~; les huit pattes de l'araignée y échouent. L'araignée
appartient à un autre groupe, quelle que soit sa ressemblance au
premier coup d'\omterm{def:g1:the-five-senses:sense}{œil}.
\end{solution}

\begin{solution}{exo:g3:sorting-living-things:7}
Pâquerette~: \omterm{def:g1:plants-around-us:parts}{fleurs} et \omterm{def:g2:from-seed-to-plant:seed}{graines}, \omterm{def:g1:plants-around-us:parts}{tige} tendre et verte. Mousse~:
jamais de \omterm{def:g1:plants-around-us:parts}{fleurs}, pas de vraie \omterm{def:g1:plants-around-us:parts}{tige}. Pin~: un arbre (tronc ligneux),
des \omterm{def:g1:plants-around-us:parts}{feuilles} en aiguilles, pas de \omterm{def:g1:plants-around-us:parts}{fleurs} mais des cônes à \omterm{def:g2:from-seed-to-plant:seed}{graines}.
Rosier~: \omterm{def:g1:plants-around-us:parts}{fleurs} et \omterm{def:g2:from-seed-to-plant:seed}{graines}, \omterm{def:g1:plants-around-us:parts}{tiges} ligneuses.
\end{solution}

\begin{solution}{exo:g3:sorting-living-things:8}
Réponse libre. Dans la plupart des rues, les \omterm{prop:g3:sorting-living-things:groups}{insectes} l'emportent de
loin --- fourmis, mouches, abeilles, scarabées, papillons --- les
\omterm{prop:g3:sorting-living-things:groups}{oiseaux} venant en second.
\end{solution}

\begin{solution}{exo:g3:sorting-living-things:9}
Celui de son amie. «~Compagnie ou sauvage~» trie selon la façon dont
les \omterm{def:g1:animals-around-us:animal}{animaux} vivent avec les gens, ce qui peut changer --- un chat
errant est le même \omterm{def:g1:animals-around-us:animal}{animal} dedans et dehors. «~\omterm{prop:g3:sorting-living-things:groups}{Mammifères} ou
\omterm{prop:g3:sorting-living-things:groups}{oiseaux}~» trie selon des caractères du corps qui restent toute la
vie~: le genre de \omterm{def:g3:sorting-living-things:criterion}{critère} du biologiste.
\end{solution}

\begin{solution}{exo:g3:sorting-living-things:10}
Étape 1~: le revêtement --- des \omterm{def:g1:animals-around-us:covering}{plumes} (serrées), un bec, deux
pattes, des ailes (en nageoires). Étape 2~: les petits sortent
d'\omterm{def:g2:animals-grow-up:born}{œufs}. Étape 3~: les caractères correspondent aux \omterm{prop:g3:sorting-living-things:groups}{oiseaux}. Étape
4~: les impressions «~nage comme un poisson, ne vole pas~»
contredisent --- et perdent. Le manchot est un oiseau.
\end{solution}
